\documentclass[11pt]{article}
\usepackage[a4paper,margin=27mm]{geometry}
\usepackage[T1]{fontenc}
\usepackage{lmodern,amsmath,amssymb,booktabs}
\usepackage[colorlinks=true,linkcolor=blue,citecolor=blue,urlcolor=blue]{hyperref}
\newcommand{\ket}[1]{|#1\rangle}
\newcommand{\Tr}{\operatorname{Tr}}
\newcommand{\dd}{\,\mathrm{d}}
\title{A balanced beam-splitter counterexample to the\\ low-order Wigner--R\'enyi vacuum bound}
\author{Xiaojian Zeng\\
\small PBC School of Finance, Tsinghua University, 100190\\
\small \texttt{xiaojianzeng@foxmail.com}}
\date{}
\begin{document}
\maketitle
\begin{abstract}
The vacuum lower bound for Wigner--R\'enyi entropy holds for balanced
beam-splitter states at orders $\alpha\geq 1/2$. We give a computer-assisted
counterexample to its proposed extension to every $0<\alpha<1/2$.
Both inputs are independently prepared in
$(20\ket{0}+\ket{2})/\sqrt{401}$, with mean photon number $2/401$ each.
After balanced mixing and discarding one output, the retained Wigner function
is exactly a Gaussian times the square of a quartic polynomial.
At $\alpha=1/4$, a directed full-plane integration bound gives
$\int W^{1/4}<9.43718$, whereas the vacuum integral is greater than $9.43892$.
The calculation treats polynomial zeros by an absolute-value chord and controls
the entire Gaussian tail analytically. It uses exact physical reconstruction
and two implementations of outward arithmetic. The result concerns one fixed
order and one finite-support product input; it does not identify an optimum
or a sharp boundary in R\'enyi order.
\end{abstract}

\section{Question and contribution}
For a state with a nonnegative Wigner function, differential entropy measures
uncertainty over phase space. Van Herstraeten and Cerf proved a vacuum
Wigner--R\'enyi entropy bound for the reduced outputs of balanced beam splitters
with separable inputs at $\alpha\geq 1/2$~\cite{interference}.
Their discussion proposes extending it to smaller positive orders.
The physical restriction matters: general Wigner-positive states admit
sub-vacuum entropy~\cite{general,he}, while the established beam-splitter
Shannon bound remains valid. Purity-based sufficient conditions for the
Shannon bound~\cite{purity} address a different regime.

Here we test the proposed low-order extension inside the narrower class of
independently prepared product inputs. The witness is an exact non-Gaussian
ket of support $\{0,2\}$ in each input. We reconstruct its reduced state
and Wigner function algebraically, then give a full-plane upper enclosure
at $\alpha=1/4$ strictly below vacuum. No finite-box replacement of the entropy
or perturbative approximation is used.

The counterexamples of Ref.~\cite{general} are explicitly outside the
beam-splitter class; its two-level output-sector discussion concerns Shannon
entropy. Our output has a nonzero four-photon component. Reference~\cite{he}
also treats general Wigner-positive states and loss stability, rather than this
fixed low-order, two-non-Gaussian-input witness.
The contribution is therefore a restricted-channel counterexample, not
another disproof of the general Shannon conjecture.

\section{Conventions and exact physical state}
We use dimensionless canonical quadratures, $[x,p]=i$, with $\hbar=1$,
$a=(x+ip)/\sqrt{2}$ and natural logarithms. Our phase-plane measure is
$\dd x\dd p$, and $R=x^2+p^2$. The vacuum Wigner function is
\begin{equation}
 W_0(x,p)=\pi^{-1}e^{-R}.
\end{equation}
For $\alpha>0$, $\alpha\neq1$, define
\begin{equation}
 I_\alpha(W)=\int_{\mathbb R^2}W^\alpha\dd x\dd p,\qquad
 h_\alpha(W)=\frac{\ln I_\alpha(W)}{1-\alpha},
 \label{eq:entropy}
\end{equation}
with $0^\alpha=0$. The vacuum satisfies
\begin{equation}
 I_\alpha(W_0)=\frac{\pi^{1-\alpha}}{\alpha},\qquad
 h_\alpha(W_0)=\ln\pi+\frac{\ln\alpha}{\alpha-1}.
 \label{eq:vac}
\end{equation}

Prepare the two modes independently in
\begin{equation}
 \ket{\psi}=\frac{20\ket{0}+\ket{2}}{\sqrt{401}}.
 \label{eq:input}
\end{equation}
Its normalization is exact and $\langle a^\dagger a\rangle=2/401$.
The balanced lossless mixer acts on creation operators by
$a^\dagger\mapsto(a^\dagger+b^\dagger)/\sqrt{2}$ and
$b^\dagger\mapsto(a^\dagger-b^\dagger)/\sqrt{2}$.
The convention differing by a sign of the second input gives the same
state here because that input has even parity. The joint output is
\begin{equation}
 \ket{\Psi}=\frac{1}{401}\left[
 400\ket{00}+20\ket{20}+20\ket{02}
 +\frac{\sqrt{6}}{4}(\ket{40}+\ket{04})-\frac12\ket{22}
 \right].
 \label{eq:joint}
\end{equation}
Tracing out the second mode gives a density matrix on $\{0,2,4\}$,
\begin{equation}
 \rho=\frac{1}{401^2}
 \begin{pmatrix}
 160400+3/8 & 7990 & 100\sqrt{6}\\
 7990 & 1601/4 & 5\sqrt{6}\\
 100\sqrt{6} & 5\sqrt{6} & 3/8
 \end{pmatrix}.
 \label{eq:rho}
\end{equation}
It is positive semidefinite as the Gram matrix obtained from
Eq.~\eqref{eq:joint}, and its trace is one. The retained mean photon number
is $2/401$. All higher Fock components vanish exactly: this support is
part of the defined state, rather than a truncation of a Gaussian input.

The corresponding Wigner function is
\begin{equation}
 W(x,p)=\frac{e^{-R}}{\pi}Q(x,p)^2,\qquad
 Q(x,p)=\frac{401-2R+R^2/2+20\sqrt{2}(x^2-p^2)}{401}.
 \label{eq:W}
\end{equation}
Consequently $W\geq0$ everywhere, including its zeros, and
$\int W\dd x\dd p=1$. Positivity of the density operator follows from its
preparation and partial trace, not merely from positivity of this
phase-space expression.

Two algebraic constructions establish Eq.~\eqref{eq:W}.
The first transforms Eq.~\eqref{eq:rho} using Fock-basis Weyl elements.
The second starts with the Hermite wavefunction
\begin{equation}
 \psi(q)=\frac{\pi^{-1/4}e^{-q^2/2}}{\sqrt{401}}
 \left(20-\frac{\sqrt{2}}{2}+\sqrt{2}q^2\right).
\end{equation}
Balanced rotation gives the joint wavefunction
$\pi^{-1/2}e^{-(q_1^2+q_2^2)/2}P(q_1,q_2)/401$, where
\begin{equation}
 P(q_1,q_2)=
 \left[20-\frac{\sqrt{2}}2+
       \frac{\sqrt{2}}2(q_1^2+q_2^2)\right]^2
       -2q_1^2q_2^2.
\end{equation}
Partial tracing in $q_2$ and applying
\begin{equation}
 W(x,p)=\frac{1}{\pi}\int_{\mathbb R}
 e^{-2ips}\langle x+s|\rho|x-s\rangle\dd s
\end{equation}
reduces to polynomial Gaussian moments. All coefficients can be compared
exactly in $\mathbb Q(\sqrt{2})$. This also directly checks normalization.
The coordinate calculation reproduces the known vacuum and identical
one-photon input baselines; the latter gives
$e^{-R}(R-1)^2/\pi$, including its nodal circle.
The implementation details and exact outputs accompany the manuscript.

\section{Full-plane integral bound at order one quarter}
At the specified order,
\begin{equation}
 I_{1/4}(W)=\pi^{-1/4}
 \int_{\mathbb R^2}e^{-R/4}\sqrt{|Q(x,p)|}\dd x\dd p,
 \qquad I_{1/4}(W_0)=4\pi^{3/4}.
 \label{eq:quarter}
\end{equation}
Fractional powers near polynomial zeros require explicit control.
We use a cellwise upper inequality rather than estimating the error of
point quadrature.

In polar coordinates $x=\sqrt R\cos\theta$, $p=\sqrt R\sin\theta$,
\begin{equation}
 Q(R,\theta)=\frac{401-2R+R^2/2+
                       20\sqrt{2}R\cos(2\theta)}{401}.
 \label{eq:polar}
\end{equation}
Each quadrant contributes equally and $\dd x\dd p=\frac12\dd R\dd\theta$.
For a cell $C=[a,b]\times[u,v]$ in the first quadrant, let
\begin{equation}
 M_C=\int_C e^{-R/4}\dd R\dd\theta,\qquad
 J_C=\int_C e^{-R/4}Q(R,\theta)\dd R\dd\theta.
\end{equation}
Suppose an outward enclosure gives $l\leq Q\leq h$ on $C$.
An upper bound $A_C$ on its weighted absolute-$Q$ integral is
\begin{equation}
 A_C=
 \begin{cases}
 J_C,&l\geq0,\\
 -J_C,&h\leq0,\\
 \displaystyle\frac{h+l}{h-l}J_C-\frac{2hl}{h-l}M_C,
       &l<0<h.
 \end{cases}
 \label{eq:chord}
\end{equation}
The last case follows from the chord of the convex function $|q|$ on
$[l,h]$. Its denominator is strictly positive.
Cauchy--Schwarz with weight $e^{-R/4}$ then gives
\begin{equation}
 \int_C e^{-R/4}\sqrt{|Q|}\dd R\dd\theta\leq\sqrt{M_CA_C}.
 \label{eq:cell}
\end{equation}
Thus cells that straddle a zero are included without dividing by $Q$ or
assuming differentiability of its fractional power.

The cell moments are analytic. Set
\begin{align}
 m_0&=4(e^{-a/4}-e^{-b/4}),\\
 m_1&=4[(a+4)e^{-a/4}-(b+4)e^{-b/4}],\\
 m_2&=4[(a^2+8a+32)e^{-a/4}-(b^2+8b+32)e^{-b/4}].
\end{align}
Writing $\Delta\theta=v-u$ and
$c_C=[\sin(2v)-\sin(2u)]/(2\Delta\theta)$, one obtains
\begin{equation}
 M_C=\Delta\theta\,m_0,\qquad
 J_C=\frac{\Delta\theta}{401}
 [401m_0-2m_1+m_2/2+20\sqrt{2}c_Cm_1].
\end{equation}
Natural interval evaluation of Eq.~\eqref{eq:polar}, with monotonicity of
$\cos(2\theta)$ on $[0,\pi/2]$, encloses $Q$.
The interior integral through $R=100$ is at most
$2\pi^{-1/4}\sum_C\sqrt{M_CA_C}$.

The omitted plane is controlled analytically. For $R\geq100$,
\begin{equation}
 |Q|\leq1+\frac{2+20\sqrt{2}}{401}R+\frac{R^2}{802}
       <2R^2.
\end{equation}
The deliberately loose last inequality already follows for $R\geq1$.
It yields
\begin{equation}
 I_{\rm tail}\leq
 \pi^{3/4}\sqrt{2}\int_{100}^{\infty}R e^{-R/4}\dd R
 =16\pi^{3/4}\sqrt{2}\,26e^{-25}
 <1.929\times10^{-8}.
 \label{eq:tail}
\end{equation}
This bound includes all angles and all radii outside the disk.

\subsection{Outward arithmetic and results}
OpenAI Codex assisted with candidate discovery, algebraic derivation,
certificate implementation and computational checking. Separate Codex sessions
provided an additional coordinate-space reconstruction and an integer-interval
implementation. The exact algebra and both certificates are supplied for
reproduction; these AI-assisted checks do not replace human scientific review.

Uniform grids use $800\times64$ or $1600\times128$ cells in
$[0,100]\times[0,\pi/2]$. Radial endpoints are exact dyadic rationals.
One implementation uses 40-digit directed Decimal arithmetic and a
slightly tighter coefficient tail bound. A separately implemented
35-digit integer interval calculation uses Eq.~\eqref{eq:tail} and
reconstructs the state in coordinates without importing the first
implementation.
Both use the cell inequality above, so this comparison checks independent
algebra and arithmetic implementations, not two unrelated integration theories.

In the integer implementation, additions and subtractions are exact;
products and divisions are enclosed by integer floor and ceiling operations.
Square roots use integer square roots with an upward correction when needed.
The Machin identity
$\pi=16\arctan(1/5)-4\arctan(1/239)$ with alternating-series remainder
encloses $\pi$. For $e^x$, $0\leq x\leq25$, terms through degree 160
are summed with a positive geometric bound on the remainder and the result
is inverted to enclose $e^{-x}$.
Sine and cosine use terms indexed $0$ through $60$ and an absolute
next-term remainder on $[0,\pi]$. The remainders lie beyond the point
where successive terms decrease. Every ensuing operation, including
the sign-changing chord denominator, is outward rounded.

\begin{table}[ht]
\centering
\small
\begin{tabular}{lrr}
\toprule
Arithmetic and grid & Full-plane upper bound & Tail upper bound\\
\midrule
Decimal, $800\times64$ & $9.437261993640953$ & $1.466\times10^{-8}$\\
Decimal, $1600\times128$ & $9.437179016856389$ & $1.466\times10^{-8}$\\
Integer, $1600\times128$ & $9.437179021478335$ & $1.929\times10^{-8}$\\
\bottomrule
\end{tabular}
\caption{Outward upper bounds on $I_{1/4}(W)$ including the entire tail.
Decimal endpoints are from \texttt{alpha-quarter-jensen-enclosures.json};
integer endpoints are from \texttt{independent-enclosure.json}.
The finer grid contains 621 sign-straddling cells.
These values bound the integral from above; they are not two-sided
enclosures of its value.}
\label{tab:bounds}
\end{table}

All three computations give a strict upper bound below the directed
vacuum lower bound. In particular,
\begin{equation}
 I_{1/4}(W)<9.43718<9.43892<4\pi^{3/4}
          =I_{1/4}(W_0).
 \label{eq:strict}
\end{equation}
The last lower bound follows from the same enclosed $\pi$; its unrounded
endpoint exceeds $9.4389219696587875$.
Since $\ln$ is increasing and $1-\alpha=3/4>0$,
Eq.~\eqref{eq:strict} proves $h_{1/4}(W)<h_{1/4}(W_0)$.
One such normalized physical product input suffices to negate a statement
quantified over every product input and every $0<\alpha<1/2$.

\section{Physical meaning and limitations}
Balanced mixing can produce a positive phase-space distribution from these
independent non-Gaussian inputs while its quarter-order differential entropy
lies below vacuum. Thus the channel's Wigner positivity and the established
bounds at larger orders do not imply the proposed bound at all positive orders.
The example has low mean occupation but contains an exact four-photon
output component; discarding that component would define a different state.
The measured quantity here is a Wigner--R\'enyi phase-space functional, not
the von Neumann entropy of the reduced density operator.

This counterexample leaves the proved $\alpha\geq1/2$ beam-splitter result,
including its Shannon limit, intact. It does not establish a violating
interval, an optimal deficit, a sharp order threshold, a universal small-input
expansion, or robustness under imperfect preparation and loss.
The full-plane comparison is computer assisted: its trust base is the
displayed inequalities, the supplied Python code and exact integer/rational
or directed Decimal operations. No formal proof-assistant verification or
experimental implementation is claimed.
Comparison with primary versions available on 12 October 2026 found no
matching restricted-class resolution; this is a bounded literature check
and not an exhaustive priority assertion.

\paragraph{Data availability and reproducibility.}
The accompanying source package includes the exact reconstruction programs,
both full-plane certificate implementations and their reference outputs.
The accompanying instructions provide commands for exact Fock
and coordinate reconstruction and both full-plane certificates in a scratch
directory.
All certificate programs require only Python's standard library.
Exact rational interval endpoints and actual execution outputs are included.
The table uses upper bounds only; no empirical search values are needed
for the conclusion.

\paragraph{Acknowledgments and AI disclosure.}
OpenAI Codex assisted with literature searching and manuscript preparation,
in addition to the research uses described above. The model used for this
manuscript preparation was \texttt{gpt-6.1-sol}; the exact model versions of
earlier calculation sessions are not fully recorded and require confirmation.
The displayed physical reconstruction and separate computational checks
describe verification performed by AI-assisted programs, not verification by
the named human author. Human authorship, contributions, review of all
AI-generated content and responsibility for the work remain to be confirmed
before submission.

\bibliographystyle{unsrt}
\bibliography{references}
\end{document}
